\section{\MakeUppercase {Background}}
\label{sec:Background}
\renewcommand{\baselinestretch}{2.0}\normalsize

This section examines how resilience has been conceptualized across disciplines and how those ideas were later taken up in management and information systems research. Resilience did not travel from one field to another as a stable concept. It developed through distinct intellectual traditions, each carrying its own assumptions about stability, adaptation, control, and system change. Those differences matter in organizational settings because they shape what resilience is taken to mean, how it is measured, and what kinds of interventions are treated as legitimate or effective. Tracing these differences helps explain why resilience has become so widely invoked in organizations while remaining difficult to define, compare, and operationalize consistently.

\subsection{How Resilience Was Shaped Across Disciplines}

\paragraph{Mathematics and Physics}

The earliest foundations of resilience emerged in mathematics and physics, where it was generally understood as the capacity of a system to absorb disturbance and return to equilibrium. Early formulations in physics and systems thinking emphasized elasticity, stability, and homeostatic feedback as defining features of resilient behavior \autocite{wiener1948cybernetics,panteli2017metrics,rankine1858manual}. In that view, resilience is mainly about rebound and temporal recovery. This perspective proved especially useful for formal and quantitative treatments of engineered systems because disturbance-response dynamics can be specified with considerable precision. Its limitation is equally clear: it leaves relatively little room for cases in which systems do not return to a prior stable state.

\paragraph{Biology and Ecology}

Biology and ecology changed the meaning of resilience more fundamentally. The emphasis moved away from equilibrium restoration and toward persistence through variability, adaptation, and regime change. Resilience no longer meant return to a prior state so much as the capacity to reorganize while preserving core functions \autocite{holling1973resilience,gunderson2002panarchy}. This perspective stresses that resilience emerges through feedbacks across ecological, social, and institutional levels \autocite{sundstrom2014transdisciplinary}. Ecological models such as the adaptive cycle illustrate transitions between stability regimes and show how diversity and redundancy can buffer systems against collapse \autocite{walker2004resilience,folke2004regime,folke2010resilience}. Early metrics such as return variance, autocorrelation, and recovery time were used to characterize resilience under uncertainty \autocite{ives1995measuring} (see Appendix \ref{appendix:A2}). This tradition matters because it treats resilience as persistence through reorganization rather than simple recovery. At the same time, it creates a problem for later organizational research: once resilience includes transformation as well as restoration, the relevant baseline, success criteria, and performance measures become harder to specify.\\

These perspectives are often treated as compatible, but they are not built on the same assumptions. Engineering resilience is oriented toward recovery to a known and valued state, whereas ecological resilience accepts irreversibility, nonlinear change, and the possibility that persistence may require transformation rather than restoration \autocite{holling1996engineering,walker2004resilience}. A further extension of this trajectory, articulated by \textcite{taleb2012antifragile}, treats some systems as not merely robust to disturbance but \emph{antifragile}: capable of gaining structurally from variability, exposure, and stress. In organizational settings the antifragile case highlights mechanisms—deliberate exposure to controlled stress, redundancy designed for option value, learning loops that compound across incidents—that can actively strengthen adaptive capacity over time, complementing both restoration- and transformation-oriented accounts. That spectrum becomes difficult to ignore in organizational research, where resilience is often invoked as if stability, adaptation, transformation, and antifragility formed a single coherent objective \autocite{woods2015four,duchek2019organizational}.

\paragraph{Neuroscience and Psychology}

Research in neuroscience and psychology extends the concept again by framing resilience as the capacity of individuals and groups to endure, adapt, and sometimes grow through adversity \autocite{luthar2000construct,kalisch2024neurobiology}. This work emphasizes mechanisms such as cognitive flexibility, learning, and emotional regulation as central to adaptive responses under stress \autocite{masten2010disasters,feder2019biology}. Empirical studies further connect resilience to stress-regulation processes and to the plasticity of adaptive capacities shaped by experience over time \autocite{mcewen2004protection,rakesh2019resilience}. Psychometric instruments such as the Wagnild and Young Scale \autocite{wagnild1993development} and the Connor-Davidson Resilience Scale \autocite{campbell2007psychometric} operationalize resilience mainly at the individual and group levels. Moving beyond strictly individual accounts, \textcite{ungar2021modeling} treats resilience as a multisystemic phenomenon distributed across personal, social, institutional, and ecological contexts. In organizational settings, especially in domains such as cybersecurity operations, healthcare, and emergency response, these insights help explain how stress, judgment, coordination, and learning shape adaptive action under pressure \autocite{norris2007community,kruk2017building,haldane2021health}. What remains less developed is the link to socio-technical organization research. Individual resilience has been studied extensively, but far less is known about how these adaptive capacities scale to collective organizational processes or interact with technological infrastructures.

\paragraph{Institutional Economics}

A major exception to individual-centered accounts appears in institutional economics, which places the collective and governance dimensions of resilience at the center. In her foundational analysis of common-pool resources, \textcite{ostrom1990governing} showed how communities operating under uncertain and demanding conditions can sustain shared resources through self-organized governance arrangements grounded in trust, reciprocity, monitoring, and graduated sanctions. In this tradition, resilience is not merely recovery or endurance. It is the capacity of a socio-economic system to sustain cooperation under uncertainty, external pressure, and resource constraints. Research on social-ecological systems similarly shows that long-term resilience depends on adaptive rule-making, polycentric governance, and alignment between incentives, actors, and shared resource systems \autocite{ostrom1990governing,lebel2006governance}. This tradition is important because it treats resilience as an emergent outcome of institutional design and collective action rather than as an individual trait or a purely technical property. What it leaves more open is how these governance principles should be translated into contemporary organizations whose coordination increasingly depends on digital infrastructures and interorganizational dependencies.

\paragraph{Cybernetics and Engineering}

Control theory and engineering reintroduced resilience into applied domains such as infrastructure, safety, and risk management. Here resilience was typically defined as the ability of a system to recover functionality after disruption, often measured through recovery time, reliability, and functional uptime \autocite{bruneau2003framework}. A foundational engineering contribution distinguishes resilience as recovery speed, reliability as frequency of success, and vulnerability as severity of failure \autocite{hashimoto1982reliability}. This equilibrium-oriented view remains influential where rapid restoration of critical functions is the primary objective (see Appendix \ref{appendix:A1}). In resilience engineering, \textcite{paries2017resilience} identified four practical capabilities -- responding, monitoring, learning, and anticipating -- as key levers for improving system resilience. This tradition offers strong operational clarity and a clear evaluative logic. Its limits become more visible, however, in settings marked by adversarial adaptation, non-stationary environments, and strategic contestation, such as cybersecurity and digitally networked systems \autocite{dupont2019cyberresilience}.

\paragraph{Complexity and Socio-technical Systems}

Complexity science and socio-technical systems theory shift the focus once more by emphasizing emergence, self-organization, and nonlinear interaction. In this view, resilience arises less from centralized control than from local interaction among heterogeneous actors and from feedback mechanisms operating across scales \autocite{levin1998ecosystems,krakovska2023resilience}. Modeling traditions such as agent-based simulation, network analysis, and game-theoretic approaches help explain how system-level resilience can emerge from distributed behavior \autocite{albert2000error,myerson2013game,gao2016universal}. Research on network robustness, for example, shows that scale-free systems may be resilient to random failures yet highly vulnerable to targeted attacks \autocite{albert2000error}. Cross-scale work on infrastructure further shows that resilience must be understood across spatial, temporal, and governance levels because disruptions frequently cascade across them \autocite{moreno2018womens}. Social-ecological systems research complements this perspective by stressing the role of institutions, collective learning, adaptive governance, and actor-resource-rule configurations in sustaining resilience \autocite{berkes1998linking,anderies2004framework,adger2005socialecological}. The structure of social networks matters as well because it conditions information flows, coordination capacity, and collective action \autocite{bodin2009role}. More recent work also makes clear that resilience strategies entail political, social, and economic trade-offs \autocite{sornette2013exploring,aase2020resilience,leichenko2024climate}. This perspective is especially useful for organizational and digital contexts because it foregrounds distributed adaptation and interdependence. Its drawback is different: it is often easier to identify resilience as an emergent system-level property than to translate that property into clear organizational responsibilities, interventions, or metrics.\\

These traditions do not converge on a single understanding of resilience. Some define it as recovery to equilibrium, others as persistence through reorganization, others as a function of cognition, governance, or distributed interaction, and Taleb's antifragility extends the trajectory further by treating exposure to disorder as a potential source of structural gain rather than as a property to be neutralized \autocite{taleb2012antifragile}. Resilience therefore entered organizational research not as a unified construct but as a composite label carrying partly incompatible assumptions about what should persist, what may change, who adapts, and how adaptation should be evaluated. Those tensions are central to understanding why resilience became attractive in organizational discourse while remaining difficult to define and operationalize consistently.

\subsection{Resilience in Organizations}

Organizational research brought resilience much closer to managerial practice. In this literature, resilience is no longer treated only as a broad system property. It is translated into questions of capability, design, governance, and adaptation. That translation was productive, but also selective. Management and organization scholars drew on several resilience traditions at once without fully resolving their underlying differences, especially over what resilience consists of, how it should be measured, and whether it refers mainly to stability, adaptation, or transformation.

\subsubsection{The Construction of Resilience in Organizations}

A large part of the organizational literature explains resilience through two broad sets of mechanisms: structural arrangements and cultural practices. Together, these strands describe how organizations prepare for disruption, respond under pressure, and learn over time \autocite{bhamra2011resilience,weick2011managing,duchek2019organizational}.\\

Structural arrangements matter because they shape how control, information, and resources move when conditions become unstable. Decentralized structures can support quicker local responses when disruptions unfold unevenly across units or time horizons \autocite{boin2013resilient}. Modularity can limit cascading failure by isolating breakdowns and preserving partial system functionality. Ambidextrous forms of organization can also help firms maintain efficiency while preserving enough flexibility to innovate under pressure \autocite{bhamra2011resilience}. These features are especially relevant in digital and cybersecurity settings, where overly centralized coordination can create bottlenecks and increase vulnerability.\\

Cultural practices matter no less. Rehearsals, simulation-based training, and pre-mortem analysis help build shared mental models, reduce uncertainty, and strengthen trust under stress \autocite{westrum2004typology,klein2008performing,weick2011managing}. Practices such as improvisation, learning orientation, and psychological safety further support teams facing novel or high-pressure situations \autocite{lanzara1999transient,vogus2007organizational}. In fast-moving threat environments, these practices make it possible to react and learn before formal control systems can be revised \autocite{westrum2004typology,duchek2019organizational}.\\

One of the main contributions of this literature is that it makes resilience organizationally interpretable. Resilience is no longer treated only as a system-level outcome but as a property that can be shaped through deliberate design choices, organizational routines, and cultural conditions. The remaining task is to identify, more precisely, the pathways through which these structural and cultural enablers actually produce adaptive behavior, and the indicators that would confirm whether such behavior is taking place.

\subsubsection{Resilience as a Dynamic Managerial Capability}

A second major line of work treats resilience as a dynamic managerial capability. From this perspective, resilience is less a reactive outcome than a path-dependent capacity formed through routines, governance arrangements, resource deployment, and strategic intent \autocite{lengnick2011developing,boin2013resilient}. This marks a shift away from narrow views of robustness or risk avoidance. The emphasis falls instead on an organization's ability to interpret uncertainty, reconfigure resources, and continue functioning under pressure.\\

Recent work explicitly frames organizational resilience as a form of dynamic capability that allows firms to sense disruption, respond effectively, and in some cases derive advantage from it \autocite{teece2007explicating,duchek2019organizational}. In that view, resilience depends on microfoundations such as routines, heuristics, and cognitive schemas that shape how actors interpret events and act under uncertainty \autocite{hepfer2022heterogeneity}. Resilience therefore cannot be reduced to resistance or continuity alone. It requires an adaptive orientation embedded in both operational and strategic action.\\

In practice, that capability is usually supported by combinations of managerial levers. Slack resources and redundancy can provide buffering capacity \autocite{barasa2018what}. Ambidextrous arrangements can support the simultaneous pursuit of stability and innovation \autocite{bhamra2011resilience}. Resilience audits and readiness diagnostics aim to assess whether organizations can absorb and adapt to shocks \autocite{burnard2011organisational,calder2021iso}. This stream is useful because it translates resilience into a language managers can act on. What it still does not resolve is a basic problem: resilience is often described as a desirable capability before there is clear agreement on what observable performance would actually demonstrate that the capability exists.

\subsubsection{Cybersecurity Resilience as a Strategic and Institutional Imperative}

Cybersecurity is a particularly revealing setting because many of the broader organizational tensions around resilience become more visible there. Cyber risks are marked by uncertainty, rapid change, and interdependent vulnerabilities that cut across technical, human, and organizational layers \autocite{linkov2013resilience}. They are also often asymmetric: threats may emerge quickly, evolve during response, and remain only partly visible until damage is already underway. Under such conditions, compliance-oriented and control-oriented approaches are not sufficient on their own \autocite{dupont2019cyberresilience}. This is one reason cybersecurity is now increasingly treated not only as a control problem, but as a resilience problem \autocite{linkov2013resilience,dupont2019cyberresilience}.\\

The cyber-resilience literature generally presents resilience as an integrated capability combining technological safeguards, strategic foresight, organizational adaptability, and institutional alignment \autocite{bjorck2015cyber,lostri2022shared}. Standards such as ISO 22301 emphasize the maintenance of core functions under disruption \autocite{calder2021iso}, while other work places more weight on learning, reflexive governance, and ongoing adjustment \autocite{hausken2020cyber}. Regulatory pressures, including those associated with Basel III and related resilience initiatives in finance, have further reinforced resilience as a strategic imperative by tying it to legitimacy, continuity, and systemic risk management \autocite{baselcommittee2011basel,almeida2020challenges}. This has contributed to a broader shift in governance: firms are increasingly judged not only by the strength of their own security posture, but also by the extent to which they help sustain resilience across wider ecosystems \autocite{bjorck2015cyber,lostri2022shared,gillard2023efficient}.\\

At the same time, this work also clarifies where current organizational thinking can be extended. Formal controls, standards, and regulatory mandates have done important work in clarifying accountability and guiding resource allocation; their next stage of development concerns how compliance and document-centered assurance can be coupled more directly to observable adaptive capacity in practice. Initial evidence from the field, including the interviews summarized in Appendix~\ref{appendix:C}, suggests this coupling is uneven and warrants closer empirical attention. AI, predictive analytics, and scenario modeling offer promising tools to support anticipation and adaptive response \autocite{kuypers2018designing,gillard2023efficient}, and their integration into routine organizational governance is steadily progressing. Across these developments, a productive tension remains between resilience as regulatory requirement, resilience as managerial capability, and resilience as an emergent property of socio-technical coordination under pressure \autocite{boin2013resilient,linkov2013resilience,duchek2019organizational}. Reconciling these three faces of resilience is precisely what the present review aims to support.\\
\\
Together, management and cybersecurity research have made resilience visible in organizational practice by identifying structural, cultural, strategic, and governance-related conditions that support adaptive action. Building on these contributions, three open questions continue to shape the broader resilience literature: how to refine the concept itself, how to converge on what counts as observable resilient performance, and how to strengthen the link between conceptual claims and operational measures. Engaging these questions is what allows resilience, already a serious organizing principle in policy and practice, to be defined and enacted with greater consistency in digitally mediated settings.